\chapter{Введение}

\section{Язык выражений}

\begin{statement}
	$ e_1 + e_2 \sim e_2 + e_1 $
\end{statement}

\begin{proof}
	Распишем эквивалентность:
	$$ \llb e_1 + e_2 \rrb \stackrel?= \llb e_2 + e_2 \rrb $$
	(по определению эквивалентности)
	$$ \llb e_1 \rrb + \llb e_2 \rrb \stackrel?=
	\llb e_2 + \rrb + \llb e_1 \rrb $$
	(по определению $ \llb \cdot \rrb $)
	$$ \llb e_1 \rrb =
	\begin{cases}
		z_1 \\
		\bot
	\end{cases}, \quad e_2 =
	\begin{cases}
		z_2 \\
		\bot
	\end{cases} $$
	Нужно рассмотреть 4 случая, получится всё, что нужно.
\end{proof}

$ 0 \sim x - x $ "--- неверно:
$$ \llb 0 \rrb = 0, \quad \llb x - x \rrb = \llb x \rrb - \llb x \rrb = \bot - \bot = \bot $$

\subsection{Экскурс в индукцию}

Так выглядит индукция в общем виде:
$$ \texttt{p\textunderscore ind} = P_\O \to \Bigl( \forall n \quad P_n \to P_{n + 1} \Bigr) \to
\forall n \quad P_n $$

\begin{minted}{OCaml}
	let P_ind pbase pstep = fun n ->
		match n with
		| 0 -> pbase
		| n -> pstep (n - 1) (P_ind pbase pstep (n - 1))
\end{minted}

Нужно доказать \texttt{pbase} и \texttt{pstep}.

$$
\begin{matrix}
	\mathscr E = & \N \\
				 & X \\
				 & \mathscr E \otimes \mathscr E
\end{matrix} $$

Доказательство по индукции для такого дерева выглядит так:
\begin{itemize}
	\item \textbf{База}
		$$ \forall n \quad P(\circled{n}) \wedge \forall x \quad P(\circled{x}) $$
	\item \textbf{Переход}
		$$ \forall e_1, e_2 \in \mathscr E \quad P(e_1) \to P(e_2) \to P \left(
			\begin{smallmatrix}
				& \otimes & \\
				\diagup & & \diagdown \\
				e_1 & & e_2
		\end{smallmatrix} \right) $$
\end{itemize}

\begin{statement}
	$ D = \Z_\bot $
	$$ \forall e \quad \llb e \rrb \ne \bot \implies FV(e) = \O $$
\end{statement}

\begin{proof}
	По \textbf{индукции}.
	\begin{itemize}
		\item \textbf{База}
			\begin{itemize}
				\item $ e = w $
					$$ \llb n \rrb = n \ne \bot $$
					$$ FV(n) = \O $$
				\item $ e = x $
					$$ \llb x \rrb = \bot $$
					$$ FV(x) = \Set{x} $$
			\end{itemize}
		\item \textbf{Переход}
			$$ e_1, e_2 $$
			$$ \llb e_1 \rrb \ne \bot \implies FV(e_1) = \O $$
			$$ \llb e_2 \rrb \ne \bot \implies FV(e_2) = \O $$
	\end{itemize}
\end{proof}

\begin{definition}
	\emph{Структурная индукция} "--- индукция по структуре программы (выражения).
\end{definition}

\begin{notation}
	$ FV(e) $ "--- множество свободных переменных $ e $.
\end{notation}

\begin{statement}
	$ \llb e \rrb \ne \bot \stackrel?\impliedby FV(e) = \O $
\end{statement}

\begin{proof}
	Хотим доказать, что
	$$ FV(e) = \O \stackrel?\implies \llb e \rrb \ne \bot $$
	Воспользуемся \textbf{структурной индукцией}.
	\begin{itemize}
		\item \textbf{База}
			\begin{itemize}
				\item $ n $
					$$ FV(n) \ne \O, \quad \llb n \rrb = n $$
				\item $ x $
					$$ FV(x) = \Set{x} \ne \O, \quad \llb x \rrb = \bot $$
			\end{itemize}
		\item \textbf{Переход}
			$$ e_1, e_2 $$
			$$ FV(e_1) = \O \implies \llb e_1 \rrb \ne \bot $$
			$$ FV(e_2) = \O \implies \llb e_2 \rrb \ne \bot $$
			$$ (e_1 \otimes e_2) = \O = FV(e_1) \cup FV(e_2) \implies
			\begin{cases}
				FV(e_1) = \O \\
				FV(e_2) = \O
			\end{cases} $$
			$$ \llb e_1 \rrb \ne \bot, \quad \llb e_2 \rrb \ne \bot, \quad
			\llb e_1 \otimes e_2 \rrb = \llb e_1 \rrb \oplus \llb e_2 \rrb \stackrel?\ne \bot $$
			Это требует, чтобы $ \oplus $ был всюду определён, а это, очевидно, \textbf{неверно}.
	\end{itemize}
\end{proof}

$$ \llb \cdot \rrb \mathscr E \to \Z_\bot $$
Хотим выражения разделить на два сорта: те, у которых есть значение, и те, у которых нет.
Это можно сделать, просто вычислив значения, но мы так не хотим.

Таким образом, нужно отображение $ \llb \cdot \rrb_s \mathscr E \mapsto \Set{\top, \bot} $:
$$ \forall e \quad \llb e \rrb_s = \top \iff \llb e \rrb \ne \bot $$

Попробуем его определить индуктивно:
$$ \llb n \rrb_s = \top $$
$$ \llb x \rrb_s = \bot $$
$$ \llb e_1 \otimes e_2 \rrb_s =
\begin{cases}
	\bot, \quad \llb e_1 \rrb_s = \bot \vee \llb e_2 \rrb_s = \bot \\
	\top, \quad \oplus \text{ всюду определена} \\
	\bot
\end{cases} $$

Надо проверить, верно ли, что так определённая $ \llb \cdot \rrb_s $
$$ \forall e \llb e \rrb_s = \top \stackrel?\iff \llb e \rrb \ne \bot $$

Это неверно, так как, например, для деления мы берём второй случай (не всюду определено), и говорим
везде $ \bot $.
А это неправда.

Определим семантику так, чтобы можно было учитывать значения переменных.
$$ \llb \cdot \rrb : \mathscr E \mapsto ??? $$

$$ \llb n \rrb = n $$
$$ \llb x \rrb = \sigma \mapsto \sigma x $$
$$ \llb e_1 \otimes e_2 \rrb = \llb e_1 \rrb \oplus \llb e_2 \rrb $$

$$ \llb \cdot \rrb : \mathscr E \mapsto \bigl( \Sigma \to \Z \bigr), \quad
\Sigma = X \to \Z $$
$ \Sigma $ называется \emph{состоянием} (\emph{окружением}).

$$ FV(e) $$
$$ \forall e, \sigma \quad \llb e \rrb \sigma \ne \bot \implies ??? $$
Понятно, что из этого не следует, что $ FV(e) = \O $.
$$ \forall e, \sigma \quad \llb e \rrb \sigma \ne \bot \implies \bigl( \forall x \in FV(e) \quad
\sigma x \text{ определена} \bigr) $$
Доказывается это \textbf{структурной индукцией} по $ e $.

$$ \llb \cdot \rrb : \mathscr E \to \bigl( \Sigma \to \Z \bigr) $$

$$ \text{Type} = \Set{\text{Int}, \text{Bool}} $$
$$ \Gamma = X \to \text{Type}_\bot $$
Построим статическую семантику
$$ \llb \cdot \rrb_s : \mathscr E \mapsto \bigl( \Gamma \to \text{Type}_\bot \bigr) $$
$$
\begin{cases}
	\llb n \rrb_s = \gamma \to \text{Type} \\
	\llb x \rrb_s = \gamma \to \gamma x \\
	\llb e_1 \otimes e_2 \rrb_s = \gamma \to
	\begin{cases}
		\bot, \quad \llb e_1 \rrb \gamma = \bot \quad \vee \quad \llb e_2 \rrb = \bot \\
		\text{Int}, \quad \llb e_1 \rrb \gamma = \text{Int} \wedge \llb e_2 \rrb \gamma = \text{Int}
		\wegde \otimes = +, -, * \\
		\bot, \quad \llb e_1 \rrb \gamma = \text{Int} \wedge \llb e_2 \rrb \gamma \text{int} \wedge
		\otimes = / \% \\
		\text{Bool}, \quad \llb e_1 \rrb \gamma = \text{Int} \wedge \llb e_2 \rrb \gamma =
		\text{Int} \wedge \otimes = <, >, \le, \ge, =, \ne \\
		\text{Bool}, \quad \llb e_1 \gamma = \text{Bool} \wegde \llb e_2 \rrb \gamma = \text{Bool}
		\wedge \otimes = \amp, \vee
	\end{cases}
\end{cases} $$

Можно записать сокращённо, но я не успел.
